\scanpage{84}
\begin{center}
APPENDIX \deleted{\unsure{to}}TO Chapter I : three letters to Larry Breen
\end{center}

\deleted{\unsure{I}} In this appendix, I am encluding three letters to Larry Breen, dated 5.2 , 17.2 and 17-19.7. 1975. These letters are written in French, and the first two are reproduced textually, whereas the third appears here in English translation. I am grateful to Ronnie Brown who took the trouble last year to make such a translation from a hardly legible \deleted{handwri} copy of the handwritten letter to Larry Breen, with the assistance of Larry Breen himself and J.L. Loday. I am now using his translation rather than the original letter, which no printer could possibly decipher ! Also, for the second letter I am using a typed copy made in 1975 by Larry Breen (who presumably had difficulties too deciphering the handwriting), Thanks are due to him for his interest and patience (with someone like me, very unknowledgeable in standard homotopy techniques), \deleted{\unsure{as it}}\inserted{\handwritten{which}} appeared \deleted{from}\inserted{\handwritten{in}} the letters I got from him in response and his verbal explanations on related matters, \inserted{\handwritten{as well as in}} the trouble he took in \deleted{copying} retyping \inserted{\handwritten{that letter}} and sending me copies of my own letters to him, and allowing me to reproduce them \deleted{\unsure{s}}in th\deleted{\unsure{a}}is volume of P\deleted{u}\handwritten{u}rsuing Stacks. My only present contribution to this set of letters is adding a few comments (in the notes (\handwritten{22}) to (\handwritten{31})), and correcting some inaccuracies in the English translation (due mainly to my handwriting ...) Also, I skipped the beginning of the first letter, which does'nt seem to me of general relevance.

The first two letters are an attempt to explain to Larry Breen (who has a wide background in algebraic geometry and homological and homotopica\unsure{l} algebra) some of the main points of the programme I had in mind\deleted{\unsure{,|x}} around the notions of $n$-categories and $n$-stacks\deleted{,} \handwritten{(}which is what I am supposed to be pursuing now in my present work ``Pursuing Stacks''\handwritten{)}. They were written under the impetus of the \inserted{new} intuition \deleted{which} (new to me at any rate), which then had just appeared to me, namely that (non strict) $n$-groupoids should model (in a suitable sense) $n$-truncated homotopy types.\handwritten{\(\mathord{/\mkern-5mu/}\)} The third letter, \inserted{\handwritten{written} in} answer to a number of questions in Larry Breen's response to the first two, is of a wider scope. A large part of the letter outlines (very ske\deleted{\unsure{ch}}\handwritten{\unsure{tch}}yly) some main points of a duality program (including a cohomological formulation of ``geometric'' local and global classfield theory), which appears here for the first time in print, \deleted{and} going back to the end of the fifties. It gives also some hints about the need of a framework of ``tame topology'' \deleted{\unsure{and}}\unsure{sui}table for writing up a ``dév\deleted{\unsure{is}}\unsure{ii|is}sage theory'' of stratified spaces, and for \deleted{\unsure{defining}}\inserted{\unsure{r|w}orking} with étale tubular neighbourhoods, for the common purpos\unsure{e} of coming to grips \deleted{of}\inserted{\handwritten{with}} a suitable notion of ``fine homotopy type'' of a ``tame\unsure{''} topological space or \inserted{\handwritten{of}} a scheme, in terms of the \deleted{ordered set}\inserted{\handwritten{inverse system}} of ``\deleted{\unsure{h}}indexed homotopy types'' corresponding to equisingular stratifications .
% archive sheet number: 83
